\begin{abstract}
This paper investigates whether acoustic emission (AE) and passive ultrasound
(US) sensing provide complementary information about the lubrication condition of
rolling-element bearings, and what such measurements actually capture. Wideband
AE and heterodyned passive US are acquired simultaneously from a single ball
bearing on an instrumented test rig, while the operating condition is varied over
a staircase of rotational speeds repeated across a sequence of temperature
blocks, driving the lubricant through the boundary, mixed, and full-film regimes.
The ISO~281 viscosity ratio $\kappa$, computed from the measured speed and
temperature, serves as a continuous reference for oil-film adequacy, and
frequency-band features from both channels are regressed onto $\kappa$ with a
regularised linear model and a gradient-boosted tree ensemble, evaluated under an
acquisition-grouped protocol that prevents leakage of near-duplicate segments
across the training--evaluation boundary.
A two-stage decomposition shows that operating-point estimation is sufficient to
reproduce the direct $\kappa$-regression performance for every channel, with no
measurable contribution beyond it: within this dataset the acoustic
$\kappa$ regression behaves as an operating-point soft sensor. Combining the two
channels nonetheless outperforms either alone---the first quantitative
confirmation of the AE--US complementarity---because they encode partly distinct
operating-condition information, ultrasound resolving temperature more strongly
than acoustic emission. These findings clarify what acoustic $\kappa$ regression
does and does not demonstrate about lubrication sensing, and indicate that
detecting lubrication changes not caused by the operating condition remains an
open problem.
\end{abstract}
